\documentclass[aspectratio=169]{beamer}

\usepackage{fontspec}
\usepackage{polyglossia}
\usepackage{color}
\usepackage{listings}
\setdefaultlanguage{russian}
\usepackage{csquotes}
\usetheme{metropolis}
\metroset{background=light}
\metroset{sectionpage=none}
\metroset{numbering=fraction}

\definecolor{ssaublue}{RGB}{77,144,205}
\setbeamercolor{frametitle}{fg=white, bg=ssaublue}
\setbeamercolor{title separator}{fg=ssaublue}

\setmainfont[
    Path = ../fonts/elektratextpro/,
    Extension = .otf,
    UprightFont = elektratextpro,
    BoldFont = elektratextpro_bold,
    ItalicFont = elektratextpro_italic,
    BoldItalicFont = elektratextpro_bolditalic
]{elektratextpro}

\setsansfont[
    Path = ../fonts/elektratextpro/,
    Extension = .otf,
    UprightFont = elektratextpro,
    BoldFont = elektratextpro_bold,
    ItalicFont = elektratextpro_italic,
    BoldItalicFont = elektratextpro_bolditalic
]{elektratextpro}

\setmonofont[
    Path = ../fonts/elektratextpro/,
    Extension = .otf,
    UprightFont = elektratextpro,
    BoldFont = elektratextpro_bold,
    ItalicFont = elektratextpro_italic,
    BoldItalicFont = elektratextpro_bolditalic
]{elektratextpro}

\newfontfamily\cyrillicfont[
    Path = ../fonts/elektratextpro/,
    Extension = .otf,
    UprightFont = elektratextpro,
    BoldFont = elektratextpro_bold,
    ItalicFont = elektratextpro_italic,
    BoldItalicFont = elektratextpro_bolditalic
]{elektratextpro}

\newfontfamily\cyrillicfontsf[
    Path = ../fonts/elektratextpro/,
    Extension = .otf,
    UprightFont = elektratextpro,
    BoldFont = elektratextpro_bold,
    ItalicFont = elektratextpro_italic,
    BoldItalicFont = elektratextpro_bolditalic
]{elektratextpro}

\newfontfamily\cyrillicfonttt[
    Path = ../fonts/elektratextpro/,
    Extension = .otf,
    UprightFont = elektratextpro,
    BoldFont = elektratextpro_bold,
    ItalicFont = elektratextpro_italic,
    BoldItalicFont = elektratextpro_bolditalic
]{elektratextpro}

% Настройка листингов для Beamer
\lstset{
    basicstyle=\ttfamily\footnotesize,
    commentstyle=\color{gray},
    columns=fullflexible,
    keepspaces=true,
    breaklines=false,
    extendedchars=true,
    texcl=true,
    showstringspaces=false
}

\title{Git в DevOps}
\author{Правдин Иван}
\date{24.05.2025}

\begin{document}
    \maketitle

    % ====================================
    % БЛОК 1: ЗАЧЕМ GIT В DEVOPS (8-10 минут)
    % ====================================

    \begin{frame}{Какую проблему решает Git в DevOps?}
        \begin{columns}

            \begin{column}{0.48\textwidth}
                \textbf{Без системы контроля версий}
                \begin{itemize}
                    \item Код только на компьютере разработчика
                    \item Нет истории изменений
                    \item Нет координации в команде
                \end{itemize}
            \end{column}

            \begin{column}{0.48\textwidth}
                \textbf{С Git}
                \begin{itemize}
                    \item Центральное хранилище кода
                    \item Полная история изменений
                    \item Совместная работа команды
                \end{itemize}
            \end{column}

        \end{columns}

        \vspace{0.5cm}

        Git — это основа для всех DevOps практик
    \end{frame}

    \begin{frame}{Git как центр DevOps процессов}
        \textbf{Git репозиторий содержит не только код:}
        \begin{itemize}
        \item Исходный код приложения
        \item Конфигурацию инфраструктуры
        \item CI/CD пайплайны
        \item Документацию
        \end{itemize}

        \vspace{0.5cm}

        \textbf{Принцип:} всё что влияет на продукт должно быть в Git
    \end{frame}

    \begin{frame}{Git события как триггеры автоматизации}
        \textbf{DevOps процессы запускаются от Git событий:}

        \begin{itemize}
            \item Push в ветку → автоматическое тестирование
            \item Pull Request → код-ревью и проверки
            \item Создание тега → релиз новой версии
            \item Merge в main → деплой на прод
        \end{itemize}

        \vspace{0.5cm}

        \textbf{Git становится пультом управления для всех процессов}
    \end{frame}

    \begin{frame}{Зачем изучать Git в контексте DevOps?}
        \begin{enumerate}
            \item Git — основа любой CI/CD системы
            \item Понимание Git workflow критично для DevOps
            \item Git hooks и webhooks — инструменты автоматизации
        \end{enumerate}
    \end{frame}

    \begin{frame}{План лекции}
        \begin{enumerate}
            \item Git workflow в DevOps командах
            \item Ветвление и стратегии слияния
            \item Git hooks для автоматизации
            \item GitHub/GitLab как DevOps платформа
            \item Webhooks для интеграций
            \item Infrastructure as Code в Git
            \item Практика: настройка автоматизации через Git
        \end{enumerate}
    \end{frame}

    % ====================================
    % БЛОК 2: GIT WORKFLOW В DEVOPS (20-25 минут)
    % ====================================

    \begin{frame}{Git Flow vs GitHub Flow}
        \begin{columns}
            \begin{column}{0.48\textwidth}
                \textbf{Git Flow}
                \begin{itemize}
                    \item Сложная модель веток
                    \item develop, release, hotfix ветки  
                    \item Подходит для больших проектов
                    \item Медленная доставка изменений
                \end{itemize}
            \end{column}

            \begin{column}{0.48\textwidth}
                \textbf{GitHub Flow}
                \begin{itemize}
                    \item Простая модель: main + feature ветки
                    \item Быстрая доставка в продакшн
                    \item Идеально для DevOps
                    \item Требует хорошие тесты и CI/CD
                \end{itemize}
            \end{column}
        \end{columns}

        \vspace{0.5cm}

        \textbf{В DevOps предпочитают GitHub Flow}
    \end{frame}

    \begin{frame}{Feature Flags и Git}
        \textbf{Проблема:} как деплоить незавершённые функции?

        \vspace{0.5cm}

        \textbf{Решение: Feature Flags}
        \begin{itemize}
            \item Код новой функции в main ветке, но отключен
            \item Включается через конфигурацию
            \item Позволяет деплоить часто без риска
        \end{itemize}

        \vspace{0.5cm}

        \textbf{Git + Feature Flags = непрерывная доставка}
    \end{frame}

    \begin{frame}{Монорепозиторий vs Мультирепозиторий}
        \begin{columns}
            \begin{column}{0.48\textwidth}
                \textbf{Монорепозиторий}
                \begin{itemize}
                    \item Весь код в одном репозитории
                    \item Проще координация изменений
                    \item Единые CI/CD пайплайны
                    \item Используют Google, Facebook
                \end{itemize}
            \end{column}

            \begin{column}{0.48\textwidth}
                \textbf{Мультирепозиторий}
                \begin{itemize}
                    \item Каждый сервис в своём репозитории
                    \item Независимые релизы
                    \item Изоляция команд
                    \item Больше настроек CI/CD
                \end{itemize}
            \end{column}
        \end{columns}

        \vspace{0.5cm}

        \textbf{Выбор зависит от размера команды и архитектуры}
    \end{frame}

    % ====================================
    % БЛОК 3: АВТОМАТИЗАЦИЯ ЧЕРЕЗ GIT (15-20 минут)  
    % ====================================

    \begin{frame}{Git Hooks: автоматизация на стороне сервера}
        \textbf{Server-side hooks:}
        \begin{itemize}
            \item pre-receive — проверка перед принятием push
            \item post-receive — действия после push
            \item post-update — обновление веб-интерфейса
        \end{itemize}

        \vspace{0.5cm}

        \textbf{Применение в DevOps:}
        \begin{itemize}
            \item Запуск тестов при push
            \item Автоматический деплой
            \item Отправка уведомлений в чат
        \end{itemize}
    \end{frame}

    \begin{frame}{Webhooks: интеграция с внешними системами}
        \textbf{Webhook = HTTP запрос при Git событии}

        \vspace{0.5cm}

        \textbf{Типичные сценарии:}
        \begin{itemize}
            \item Push → запуск Jenkins job
            \item Pull Request → уведомление в Slack
            \item Tag создан → релиз в производство
            \item Issue создана → задача в Jira
        \end{itemize}

        \vspace{0.5cm}

        \textbf{GitHub, GitLab, Bitbucket поддерживают webhooks}
    \end{frame}

    \begin{frame}{GitHub Actions: CI/CD встроенный в Git}
        \textbf{YAML файлы в .github/workflows/}
        
        \vspace{0.5cm}
        
        Автоматически выполняют:
        \begin{itemize}
            \item Сборку и тестирование кода
            \item Линтинг и проверки безопасности  
            \item Деплой в различные окружения
            \item Создание релизов
        \end{itemize}

        \vspace{0.5cm}

        \textbf{Git событие → GitHub Action → результат}
    \end{frame}

    % ====================================
    % БЛОК 4: INFRASTRUCTURE AS CODE (10 минут)
    % ====================================

    \begin{frame}{Infrastructure as Code в Git}
        \textbf{Принцип:} инфраструктура описывается кодом и хранится в Git

        \vspace{0.5cm}

        \textbf{Инструменты:}
        \begin{itemize}
            \item Terraform — создание инфраструктуры
            \item Ansible — настройка серверов  
            \item Kubernetes YAML — оркестрация контейнеров
            \item Docker Compose — локальная разработка
        \end{itemize}

        \vspace{0.5cm}

        \textbf{Изменение инфраструктуры = Pull Request}
    \end{frame}

    \begin{frame}{GitOps: Git как источник истины}
        \textbf{GitOps = Infrastructure as Code + автоматизация}

        \vspace{0.5cm}

        \textbf{Принципы GitOps:}
        \begin{itemize}
            \item Декларативное описание системы в Git
            \item Git как единственный источник истины
            \item Автоматическая синхронизация с кластером
            \item Откат изменений через Git revert
        \end{itemize}

        \vspace{0.5cm}

        \textbf{Инструменты: ArgoCD, Flux, Jenkins X}
    \end{frame}

    % ====================================
    % ЗАКЛЮЧЕНИЕ
    % ====================================

    \begin{frame}{Что запомнить}
        \begin{itemize}
            \item Git — центральный элемент DevOps процессов
            \item Простые workflow лучше сложных в DevOps
            \item Webhooks и hooks — основа автоматизации
            \item Infrastructure as Code должен быть в Git
            \item GitOps — эволюция DevOps практик
        \end{itemize}

        \vspace{1cm}

        \begin{center}
            \textbf{Вопросы?}
        \end{center}
    \end{frame}

\end{document}